\documentclass[../../main.tex]{subfiles}
\begin{document}
% SOURCE p=239; printed=224; object=ch07-czero-exercises-title; region=lower
\subsection*{习题}
% SOURCE p=239; printed=224; object=ch07-ex16-functional-measure-jordan; number=16; region=lower
\begin{exercise}\label{exercise:ch07-functional-measure-jordan}
设 $I\in C_0(X,\mathbb R)^*$，$I^+,I^-$ 是\cref{lemma:ch07-functional-jordan}证明中构造的泛函。若 $\mu$ 为与 $I$ 关联的符号 Radon 测度，则 $\mu$ 的正、负变差就是分别与 $I^+,I^-$ 关联的 Radon 测度。
\end{exercise}
% SUPPLEMENT: id=v14-ch07-ex16-solution; parent=ch07-ex16-functional-measure-jordan; kind=solution+motivation+hypotheses; origin=independent-derivation
\begin{solution}
\textbf{完整补写（重构新增）。}这里 $\mu$ 有限，因为它表示有界泛函。由\cref{proposition:ch07-variation-radon}，$\mu^+,\mu^-$ 为有限 Radon 测度。取 $\mu$ 的 Borel Hahn 分解 $X=P\sqcup N$。对 $f\in C_0(X,\mathbb R)$、$f\geqslant0$ 及任意 $0\leqslant g\leqslant f$，
\begin{align*}
I(g)=\int g\,\mathrm d\mu^+-\int g\,\mathrm d\mu^-
\leqslant\int f\,\mathrm d\mu^+,
\end{align*}
取上确界即得 $I^+(f)\leqslant\int f\,\mathrm d\mu^+$。

由于 $|\mu|$ 有限，$\chi_P\in L^1(|\mu|)$。用 $C_c$ 的 $L^1$ 稠密性取实 $u_n\in C_c(X)$ 使 $u_n\to\chi_P$ 于 $L^1(|\mu|)$，再令 $h_n=\max(0,\min(1,u_n))$。截断不增加与 $0$ 或 $1$ 的距离，故 $0\leqslant h_n\leqslant1$，$h_n\in C_c$，且 $\|h_n-\chi_P\|_{L^1(|\mu|)}\to0$。置 $g_n=fh_n$，则 $g_n\in C_0$ 且 $0\leqslant g_n\leqslant f$，并且
\begin{align*}
\left|I(g_n)-\int f\,\mathrm d\mu^+\right|
&=\left|\int f(h_n-\chi_P)\,\mathrm d\mu\right|\\*
&\leqslant\|f\|_u\|h_n-\chi_P\|_{L^1(|\mu|)}\longrightarrow0.
\end{align*}
因此 $I^+(f)\geqslant\int f\,\mathrm d\mu^+$，两边相等。每个实 $f$ 是其非负部分之差，线性性使等式扩展到整个实 $C_0$。最后 $I^-=I^+-I$ 给出 $I^-(f)=\int f\,\mathrm d\mu^-$，由表示唯一性完成识别。
\end{solution}
\begin{remark}
\textbf{动机与条件（重构新增）。}抽象的正锥上确界恰好选择 Hahn 正集上的质量；连续函数无法直接取 $f\chi_P$，所以要用有限 Radon 测度下的 $L^1$ 逼近恢复它。$I^\pm$ 必须是指定上确界构造的分量，任意正泛函差的分解可以同时增加同一个正泛函，不具有此唯一性。
\end{remark}
% END SUPPLEMENT: id=v14-ch07-ex16-solution
% SOURCE p=239; printed=224; object=ch07-ex17-positive-czero-automatic-bound; number=17; region=lower
\begin{exercise}\label{exercise:ch07-positive-czero-automatic-bound}
若 $\mu$ 为 $X$ 上满足 $\mu(X)=\infty$ 的正 Radon 测度，则存在 $f\in C_0(X)$ 使 $\int f\,\mathrm d\mu=\infty$。从而 $C_0(X)$ 上每个正线性泛函都有界。
\end{exercise}
% SUPPLEMENT: id=v14-ch07-ex17-solution; parent=ch07-ex17-positive-czero-automatic-bound; kind=solution+motivation+hypotheses; origin=independent-derivation
\begin{solution}
\textbf{完整补写（重构新增）。}开集 $X$ 的质量公式使我们对每个 $n\geqslant1$ 都能选择 $g_n\in C_c(X)$，满足 $0\leqslant g_n\leqslant1$ 及 $\int g_n\,\mathrm d\mu>4^n$。级数 $f=\sum_{n=1}^\infty2^{-n}g_n$ 一致收敛；其部分和属于 $C_0$，而 $C_0$ 一致闭，故 $f\in C_0$，且 $f\geqslant0$。对每个 $n$，
\begin{align*}
\int f\,\mathrm d\mu\geqslant2^{-n}\int g_n\,\mathrm d\mu>2^n,
\end{align*}
所以积分为无穷大。

现设 $I$ 为 $C_0(X)$ 上的正线性泛函，尚不假设连续。它在 $C_c$ 上的限制由正 Riesz 定理表示成正 Radon 测度 $\mu$。若 $\mu(X)=\infty$，对上面构造的 $f$ 及部分和 $s_N$ 有 $0\leqslant s_N\leqslant f$，从而
\begin{align*}
I(f)\geqslant I(s_N)=\int s_N\,\mathrm d\mu
>\sum_{n=1}^N2^n,
\end{align*}
与 $I(f)$ 为有限数矛盾，故 $\mu(X)<\infty$。

还须证明原 $I$ 在 $C_0$ 上就是这个有限测度积分，不能先行假设它连续。对非负 $f\in C_0$ 及 $\varepsilon>0$，令 $f_\varepsilon=(f-\varepsilon)^+$，其支集包含于紧集 $\{f\geqslant\varepsilon\}$，故 $f_\varepsilon\in C_c$。逐点有
\begin{align*}
0\leqslant f-f_\varepsilon=\min(f,\varepsilon)
\leqslant\sqrt\varepsilon\sqrt f.
\end{align*}
这里 $\sqrt f\in C_0$，正性给 $0\leqslant I(f)-I(f_\varepsilon)\leqslant\sqrt\varepsilon I(\sqrt f)\to0$。取 $\varepsilon=1/n$，单调收敛又给 $\int f_{1/n}\,\mathrm d\mu\to\int f\,\mathrm d\mu$。因此 $I(f)=\int f\,\mathrm d\mu$。实函数分正负部分、复函数再分实虚部即可推广到全部 $C_0$。于是 $|I(f)|\leqslant\mu(X)\|f\|_u$，证明有界，且前面的范数公式给 $\|I\|=\mu(X)$。
\end{solution}
\begin{remark}
\textbf{动机与条件（重构新增）。}小的一致幅度与越来越大的积分相乘，可以把无穷总质量汇集到一个仍在无穷远消失的函数中。正性排除了抵消，使部分和受 $I(f)$ 控制；定义在整个 $C_0$ 上保证级数极限也能代入 $I$。单有 $C_c$ 上的正性只给局部界，例如 Lebesgue 积分在一致范数下并不有界。最后的截断论证专门弥补“稠密子空间相等”在未知连续时不能推出处处相等的逻辑缺口。
\end{remark}
% END SUPPLEMENT: id=v14-ch07-ex17-solution
% SOURCE p=239; printed=224; object=ch07-ex18-radon-lebesgue-decomposition; number=18; region=bottom
\begin{exercise}\label{exercise:ch07-radon-lebesgue-decomposition}
若 $\mu$ 为 $X$ 上的 $\sigma$ 有限 Radon 测度，$\nu\in M(X)$，且 $\nu=\nu_1+\nu_2$ 是 $\nu$ 关于 $\mu$ 的 Lebesgue 分解，则 $\nu_1,\nu_2$ 都是 Radon 测度。（使用\cref{exercise:ch07-integrable-weight-radon}。）
\end{exercise}
% SUPPLEMENT: id=v14-ch07-ex18-solution; parent=ch07-ex18-radon-lebesgue-decomposition; kind=solution+motivation+hypotheses; origin=independent-derivation
\begin{solution}
\textbf{完整补写（重构新增）。}按 Lebesgue 分解的约定，$\nu_1\ll\mu$、$\nu_2\perp\mu$，且由复 Radon--Nikodym 定理有 $\nu_1=f\mu$，其中 $f\in L^1(\mu)$。因此 $|\nu_1|=|f|\mu$，由\cref{exercise:ch07-integrable-weight-radon}为有限 Radon 测度，再由\cref{proposition:ch07-variation-radon}知 $\nu_1\in M(X)$。同一命题中的向量空间性质给 $\nu_2=\nu-\nu_1\in M(X)$，结论成立。
\end{solution}
\begin{remark}
\textbf{动机与条件（重构新增）。}测度论上的绝对连续与奇异分解仍留在拓扑正则的测度类内。$\sigma$ 有限性允许直接使用第三章的分解定理，$\nu$ 的有限全变差保证所得密度可积；Radon 条件用于把密度乘出的测度以及差测度保持在 $M(X)$ 内。
\end{remark}
% END SUPPLEMENT: id=v14-ch07-ex18-solution
% SOURCE p=240; printed=225; object=ch07-ex19-algebra-dual; number=19; region=top
\begin{exercise}\label{exercise:ch07-algebra-dual}
设 $X$ 为完全正则空间，$\mathcal A$ 为 $BC(X)$ 的完全正则子代数（见\cref{exercise:ch04-compactification-algebras}）。把 $\mathcal A^*$ 描述成一个测度空间。
\end{exercise}
% SUPPLEMENT: id=v14-ch07-ex19-solution; parent=ch07-ex19-algebra-dual; kind=solution+motivation+hypotheses+source-clarification; origin=independent-derivation
\begin{solution}
\textbf{按原定义及其修正版分别解答（重构新增）。}先严格沿用原题的复代数定义，它未要求对共轭封闭。每个 $f\in BC(X)$ 都唯一延拓为 $\widehat f\in C(\beta X)$，且由于 $X$ 在 $\beta X$ 稠密，这个延拓给出等距同构 $BC(X)\cong C(\beta X)$。令
\begin{align*}
\mathcal A^\perp=\left\{\lambda\in M(\beta X):
\int\widehat f\,\mathrm d\lambda=0\text{ 对所有 }f\in\mathcal A\right\}.
\end{align*}
这是闭线性子空间，因为每个积分映射都在全变差范数下连续。定义 $R:M(\beta X)\to\mathcal A^*$ 为 $R\lambda(f)=\int\widehat f\,\mathrm d\lambda$，则 $\|R\lambda\|\leqslant\|\lambda\|$，且核恰为 $\mathcal A^\perp$。

任取 $L\in\mathcal A^*$。Hahn--Banach 定理把它等范数延拓到 $C(\beta X)$；紧空间 Riesz 定理给 $\lambda\in M(\beta X)$ 表示该延拓，且 $\|\lambda\|=\|L\|$。所以 $R$ 满射，并诱导等距同构
\begin{align*}
\mathcal A^*\cong M(\beta X)/\mathcal A^\perp,\qquad
\|[\lambda]\|=\inf_{\eta\in\mathcal A^\perp}\|\lambda+\eta\|.
\end{align*}
确实，收缩性给 $\|R\lambda\|\leqslant\|[\lambda]\|$；对 $R\lambda$ 用刚才的等范数延拓得到同一陪集中的代表，给出反向不等式，且下确界可达到。这是原定义下总是正确的测度商空间描述。

若按\cref{exercise:ch04-compactification-algebras}中已证明的订正版，另要求 $\mathcal A$ 对共轭封闭，则它恰为某个 Hausdorff 紧化 $(Y,e)$ 上 $C(Y)$ 的限制代数。限制映射 $C(Y)\to\mathcal A$ 是等距满射，所以每个 $L\in\mathcal A^*$ 唯一具有形式
\begin{align*}
L(f)=\int_Y\widetilde f\,\mathrm d\lambda,\qquad\lambda\in M(Y),
\end{align*}
其中 $\widetilde f$ 为 $f$ 在 $Y$ 上的唯一连续延拓，且 $\|L\|=\|\lambda\|$。于是该修正版中可以直接写 $\mathcal A^*\cong M(Y)$，无需取商。
\end{solution}
\begin{remark}
\textbf{动机与条件（重构新增）。}代数只保留一部分测试函数，因此不同测度可能给出相同读数，商掉它们不可区分的差便得到对偶。完全正则性保证 Stone--Čech 紧化及延拓可用；原定义中的闭性使 $\mathcal A$ 为 Banach 空间。若要把全部测试函数直接识别成一个紧化上的 $C(Y)$，共轭封闭不可悄悄省略；第四章的反例已说明原定义不足以保证这种表示。对偶商空间论证本身甚至适用于 $BC(X)$ 的任意闭线性子空间。
\end{remark}
% END SUPPLEMENT: id=v14-ch07-ex19-solution
\subfile{section-03-exercises-convergence.tex}
\end{document}
